\chapter{Recovering the latents: critical velocity and anaerobic reserve}
\label{ch:energetics}

\section{The critical-power model of a running horse}

The drift $v_i$ and the reserve $\Wp_i$ of Chapter~\ref{ch:gen} are not free numbers; they are the
two parameters of a \emph{critical-power} physiology, the standard mechanistic model of endurance
performance imported from human exercise science. A horse is idealised as possessing
\begin{itemize}
  \item a \textbf{critical velocity} $\vc$ (m/s): the fastest pace it can hold purely
        \emph{aerobically}, i.e.\ indefinitely on the timescale of a race; and
  \item a finite \textbf{anaerobic reserve} $\Wp$: a fixed quantity of work it can spend running
        \emph{above} $\vc$, after which it must drop back to (or below) $\vc$.
\end{itemize}
The defining relation of the critical-power model is hyperbolic: the distance $L$ a horse can
cover in time $T$ satisfies $L = \vc\,T + \Wp$, i.e.\ it sustains $\vc$ for the whole trip and the
reserve $\Wp$ buys an extra block of distance. Solving for time gives \emph{Hill's
critical-velocity law},
\begin{equation}
  T \;=\; \frac{L - \Wp}{\vc},
  \label{eq:hill}
\end{equation}
which is exactly the deterministic skeleton behind the production mean
\eqref{eq:mu-wprime}: dividing numerator and denominator by $L$,
$T = \frac{L}{\vc}\bigl(1 - \Wp/L\bigr) \approx \frac{L}{\vc(1+\Wp/L)^{-1}}$, and the production
code carries the reserve as the multiplicative late boost $(1+\Wp/L)$ on the drift, which is the
same first-order effect with a numerically gentler, strictly-positive form. The units used
throughout are metres, seconds, and m/s; $\Wp$ is carried as the work proxy
$\int (v-\vc)^{+}\,dt$ with units $\mathrm{m^2/s}$, which is monotone in the true energetic
reserve and suffices for a relative model.

\section{The inverse problem: $(\vc,\Wp)$ from one run's sectionals}

HKJC publishes \emph{sectional} split times: a race is divided into segments, and the time to
cover each segment is recorded. By HKJC's verified convention the final segments are each
$\SI{400}{m}$ and the first segment carries the variable remainder, so for a race of length $L$
split into $n$ sections the per-section distances are
\begin{equation}
  d_1 = L - 400(n-1), \qquad d_2=\dots=d_n = 400,
  \label{eq:sectiondist}
\end{equation}
and the per-section velocities are $v_k = d_k / t_k$ from the published split times $t_k$. The
velocity profile $(v_1,\dots,v_n)$ is the raw kinematic signature of a single run, and from it we
recover that run's $(\vc,\Wp)$ by inverting the critical-power picture.

\paragraph{Estimating $\vc$.} The critical velocity is the speed the horse \emph{settles to}
once the reserve is spent --- the late-race sustainable pace. It is estimated robustly as the
slower of the final two section velocities,
\begin{equation}
  \widehat{\vc} \;=\; \min(v_{n-1},\, v_{n}),
  \label{eq:vc-est}
\end{equation}
the $\min$ guarding against a final-section sprint (a horse with reserve still in hand kicks on
the line, which would over-read $\vc$). Typical recovered values sit around
$\widehat{\vc}\approx \SI{16.8}{m/s}$.

\paragraph{Estimating $\Wp$.} The reserve is the work run \emph{above} that critical velocity,
accumulated over the sections where the horse exceeded it:
\begin{equation}
  \widehat{\Wp} \;=\; \sum_{k=1}^{n} \bigl(v_k - \widehat{\vc}\bigr)^{+}\, t_k,
  \label{eq:wp-est}
\end{equation}
where $(\cdot)^{+}=\max(\cdot,0)$. Two auxiliary descriptors of the same profile are retained for
the trajectory-shape features of Chapter~\ref{ch:features}: the peak section velocity
$v_{\text{peak}}=\max_k v_k$ (the anaerobic ceiling) and the fade
$v_{\text{peak}}/v_n - 1$ (how far the horse dropped from its peak by the line).

\section{Why per-run estimates are noisy --- and must be pooled}

The estimates \eqref{eq:vc-est}--\eqref{eq:wp-est} are \emph{per-run} and individually noisy, for
reasons that are physical rather than statistical artefacts:
\begin{itemize}
  \item an easy winner is eased down through the line, depressing $v_{n}$ and hence under-reading
        $\widehat{\vc}$;
  \item a horse caught in traffic runs a distorted profile that reflects the race shape, not its
        capacity;
  \item the pace of the race --- which the horse does not choose alone --- sets how much of the
        reserve is visible.
\end{itemize}
The latent $(\vc,\Wp)$ are properties of the \emph{horse}, approximately stable across its
career; any single run is a noisy glimpse of them through the lens of that day's tactics. This is
precisely the structure a hierarchical model is built for, and it motivates
Chapter~\ref{ch:pop}: rather than trust one run, we pool a horse's runs toward a population that
is itself conditioned on the horse's class and trip, with the degree of pooling set by how much
evidence the horse has provided. The per-run estimator of this chapter is the \emph{measurement};
the population model of the next chapter is the \emph{estimator of the latent} from those
measurements.

\begin{remark}[A note on what $\vc$ ``is'' in the production model]
The production drift is built around $\vc$ as the organising latent, but $\vc$ should be read as
a \emph{calibrated relative ability in velocity units}, not as a literal physiological constant
measured to three significant figures. Its value is anchored by the sectional estimator above, but
it is then shrunk, rating-conditioned, and nudged by the ability channels of
Chapters~\ref{ch:features}--\ref{ch:competition}, all of which live in the same m/s units so that
they add coherently. The physiology fixes the \emph{structure} (a sustainable ceiling plus a
finite reserve, distance-aware through $\Wp/L$); the calibration fixes the \emph{numbers}.
\end{remark}

\section{A generative race model from $(\vc,\Wp)$ alone}

Before the full stochastic apparatus of Chapter~\ref{ch:gen}, it is worth recording the simplest
generative model the energetics already support, because it shows the mechanism cleanly. Given
per-horse $(\vc_i,\Wp_i)$ and a trip $L$, Hill's law \eqref{eq:hill} gives a deterministic
predicted time $\widehat T_i = (L-\Wp_i)/\vc_i$; adding race-day noise $\varepsilon_i$ and ranking
yields a finishing order,
\begin{equation}
  \Tfin_i \;=\; \frac{L-\Wp_i}{\vc_i} + \varepsilon_i,
  \qquad
  \text{winner} = \argmin_i \Tfin_i .
  \label{eq:naive-gen}
\end{equation}
This is dynamics, not regression: change the distance, or a horse's $(\vc_i,\Wp_i)$, and the
whole finishing-order distribution moves mechanistically --- a horse with a large reserve gains
relatively more over a sprint than over a staying trip, automatically, because $\Wp$ enters as
$\Wp/L$. The production model replaces the additive-Gaussian noise of \eqref{eq:naive-gen} with
the physically-correct first-passage (Wald) noise of Chapter~\ref{ch:gen} --- strictly positive,
right-skewed, with a single diffusion constant and the $\Var\propto L$ law --- but the skeleton is
this critical-power generator.
